% GENERATED FILE. Edit canonical lessons, then run npm run generate:books.
% canonical-source-hash: fnv1a64:82b48f5781654b80
% canonical-lessons: ZH-C483-haozhuan, ZH-C483-kangfu, ZH-C483-yingyang, ZH-C483-baojian, ZH-C483-manxing

\chapter{Getting Better}
\label{ch:zh-getting-better}
\hlchaptermodality{483}

\begin{chapteropening}
\textbf{By the end of this chapter:} \emph{I can say someone's illness is improving, that someone has recovered, and talk about nourishing food, keeping well and long illnesses.}

Complete the last lesson of chapter 483: I can say someone's illness is improving, that someone has recovered, and talk about nourishing food, keeping well and long illnesses.
\end{chapteropening}

\section[hǎozhuǎn]{hǎozhuǎn — to take a turn for the better; to start to improve}

\label{lesson:ZH-C483-haozhuan}



\begin{quote}
Before the new one: say the Mandarin for to disinfect, then the Mandarin for to care for the sick.
\end{quote}

\subsection*{You'll want to know: hǎozhuǎn}
\emph{hǎozhuǎn} — \textquotedblleft{}to take a turn for the better; to start to improve\textquotedblright{}.

\begin{quote}
\emph{tā zhìliáo le yí ge yuè bìngqíng hǎozhuǎn le} — He had treatment for a month, and his illness has taken a turn for the better.
\end{quote}

Its characters come later; for now it is a word for the ear.

Two third tones meet in this word: hǎo is said with a rising tone before zhuǎn, and the spelling keeps both marks.

\begin{cousinweb}
\emph{hǎozhuǎn} is built from \zh{好} 'good' and \emph{zhuǎn} 'to turn', both inherited from Middle Chinese, \emph{xawX} and \emph{trjwenX} (Guangyun, in Baxter's transcription), which glosses the second 'to move' and 'to turn round'.
\end{cousinweb}

\begin{grammarlens}[title={What you've built}]
One more word for saying how it was.
\end{grammarlens}

\subsection*{Your turn}
\begin{itemize}
\raggedright
  \item \emph{Say it:} \emph{hǎozhuǎn}
  \item \emph{Say it:} \emph{hǎozhuǎn}, once more
  \item \emph{Say it:} \emph{tā zhìliáo le yí ge yuè bìngqíng hǎozhuǎn le} — He had treatment for a month, and his illness has taken a turn for the better.
  \item \emph{Recall:} say the Mandarin for to disinfect, then the Mandarin for to care for the sick, then say \emph{hǎozhuǎn} again
\end{itemize}

\begin{tcolorbox}[breakable,colback=teal!4,colframe=teal!35!black,title={Before you move on}]
What does hǎozhuǎn mean? (\textquotedblleft{}To take a turn for the better; to start to improve\textquotedblright{}.) Say it once more.
\end{tcolorbox}
\section[kāngfù]{kāngfù — to get well again; to convalesce}

\label{lesson:ZH-C483-kangfu}



\begin{quote}
Before the new one: say the Mandarin for to care for the sick, then the Mandarin for to start to improve.
\end{quote}

\subsection*{You'll want to know: kāngfù}
\emph{kāngfù} — \textquotedblleft{}to get well again; to convalesce\textquotedblright{}.

\begin{quote}
\emph{tā kāngfù le xià ge yuè kěyǐ shàngbān} — He has recovered and can go back to work next month.
\end{quote}

Its characters come later; for now it is a word for the ear.

\begin{cousinweb}
\emph{kāngfù} is built from \emph{kāng} 'at ease' and \emph{fù} 'to return', both inherited from Middle Chinese, \emph{khang} and \emph{bjuwk} (Guangyun, in Baxter's transcription), which glosses the first 'peaceful' and 'happy' and the second 'to return'.
\end{cousinweb}

\begin{grammarlens}[title={What you've built}]
One more word for saying how it was.
\end{grammarlens}

\subsection*{Your turn}
\begin{itemize}
\raggedright
  \item \emph{Say it:} \emph{kāngfù}
  \item \emph{Say it:} \emph{kāngfù}, once more
  \item \emph{Say it:} \emph{tā kāngfù le xià ge yuè kěyǐ shàngbān} — He has recovered and can go back to work next month.
  \item \emph{Recall:} say the Mandarin for to care for the sick, then the Mandarin for to start to improve, then say \emph{kāngfù} again
\end{itemize}

\begin{tcolorbox}[breakable,colback=teal!4,colframe=teal!35!black,title={Before you move on}]
What does kāngfù mean? (\textquotedblleft{}To get well again; to convalesce\textquotedblright{}.) Say it once more.
\end{tcolorbox}
\section[yíngyǎng]{yíngyǎng — nutrition; nourishment}

\label{lesson:ZH-C483-yingyang}



\begin{quote}
Before the new one: say the Mandarin for to start to improve, then the Mandarin for to get well again.
\end{quote}

\subsection*{You'll want to know: yíngyǎng}
\emph{yíngyǎng} — \textquotedblleft{}nutrition; nourishment\textquotedblright{}.

\begin{quote}
\emph{shuǐguǒ hěn yǒu yíngyǎng nǐ yīnggāi měi tiān chī} — Fruit is very good for you, so you should eat some every day.
\end{quote}

Its characters come later; for now it is a word for the ear.

\begin{cousinweb}
\emph{yíngyǎng} is built from \emph{yíng} 'to build' and \emph{yǎng} 'to nourish', both inherited from Middle Chinese, \emph{yweng} and \emph{yangX} (Guangyun, in Baxter's transcription), which glosses the first 'to make' and the second 'to raise' and 'to feed'.
\end{cousinweb}

\begin{grammarlens}[title={What you've built}]
One more word for saying how it was.
\end{grammarlens}

\subsection*{Your turn}
\begin{itemize}
\raggedright
  \item \emph{Say it:} \emph{yíngyǎng}
  \item \emph{Say it:} \emph{yíngyǎng}, once more
  \item \emph{Say it:} \emph{shuǐguǒ hěn yǒu yíngyǎng nǐ yīnggāi měi tiān chī} — Fruit is very good for you, so you should eat some every day.
  \item \emph{Recall:} say the Mandarin for to start to improve, then the Mandarin for to get well again, then say \emph{yíngyǎng} again
\end{itemize}

\begin{tcolorbox}[breakable,colback=teal!4,colframe=teal!35!black,title={Before you move on}]
What does yíngyǎng mean? (\textquotedblleft{}Nutrition; nourishment\textquotedblright{}.) Say it once more.
\end{tcolorbox}
\section[bǎojiàn]{bǎojiàn — health care; looking after one's health}

\label{lesson:ZH-C483-baojian}



\begin{quote}
Before the new one: say the Mandarin for to get well again, then the Mandarin for nourishment.
\end{quote}

\subsection*{You'll want to know: bǎojiàn}
\emph{bǎojiàn} — \textquotedblleft{}health care; looking after one's health\textquotedblright{}.

\begin{quote}
\emph{wǒ yéye bǐ wǒ nǎinai zhùyì bǎojiàn} — My grandfather takes more care of his health than my grandmother does.
\end{quote}

Its characters come later; for now it is a word for the ear.

\begin{cousinweb}
\emph{bǎojiàn} is built from \zh{保} 'to keep' and \emph{jiàn} 'strong', both inherited from Middle Chinese, \emph{pawX} and \emph{gjonH} (Guangyun, in Baxter's transcription), which glosses the first 'to keep safe' and the second 'strong'.
\end{cousinweb}

\begin{grammarlens}[title={What you've built}]
One more word for saying how it was.
\end{grammarlens}

\subsection*{Your turn}
\begin{itemize}
\raggedright
  \item \emph{Say it:} \emph{bǎojiàn}
  \item \emph{Say it:} \emph{bǎojiàn}, once more
  \item \emph{Say it:} \emph{wǒ yéye bǐ wǒ nǎinai zhùyì bǎojiàn} — My grandfather takes more care of his health than my grandmother does.
  \item \emph{Recall:} say the Mandarin for to get well again, then the Mandarin for nourishment, then say \emph{bǎojiàn} again
\end{itemize}

\begin{tcolorbox}[breakable,colback=teal!4,colframe=teal!35!black,title={Before you move on}]
What does bǎojiàn mean? (\textquotedblleft{}Health care; looking after one's health\textquotedblright{}.) Say it once more.
\end{tcolorbox}
\section[mànxìng]{mànxìng — chronic (of an illness)}

\label{lesson:ZH-C483-manxing}



\begin{quote}
Before the new one: say the Mandarin for nourishment, then the Mandarin for looking after one's health.
\end{quote}

\subsection*{You'll want to know: mànxìng}
\emph{mànxìng} — \textquotedblleft{}chronic (of an illness)\textquotedblright{}.

\begin{quote}
\emph{tā yǒu mànxìngbìng yào chángcháng chī yào} — He has a chronic illness and has to take medicine regularly.
\end{quote}

Its characters come later; for now it is a word for the ear.

\begin{cousinweb}
\emph{mànxìng} is built from \emph{màn} 'slow' and \zh{性} 'nature', both inherited from Middle Chinese, \emph{maenH} and \emph{sjengH} (Guangyun, in Baxter's transcription), which glosses the first 'idle' and the second 'nature and conduct'.
\end{cousinweb}

\begin{grammarlens}[title={What you've built}]
One more word for saying how it was.
\end{grammarlens}

\subsection*{Your turn}
\begin{itemize}
\raggedright
  \item \emph{Say it:} \emph{mànxìng}
  \item \emph{Say it:} \emph{mànxìng}, once more
  \item \emph{Say it:} \emph{tā yǒu mànxìngbìng yào chángcháng chī yào} — He has a chronic illness and has to take medicine regularly.
  \item \emph{Recall:} say the Mandarin for nourishment, then the Mandarin for looking after one's health, then say \emph{mànxìng} again
\end{itemize}

\begin{tcolorbox}[breakable,colback=teal!4,colframe=teal!35!black,title={Before you move on}]
What does mànxìng mean? (\textquotedblleft{}Chronic (of an illness)\textquotedblright{}.) Say it once more.
\end{tcolorbox}
